% Shared TikZ preamble for the m-format paper-computer wheels.
% Engine: xelatex. Distinct from docs/manual/pandoc/preamble.tex.
\documentclass[letterpaper]{article}
\usepackage[letterpaper, margin=0.5in, nohead, nofoot]{geometry}
\usepackage{fontspec}
\usepackage{tikz}
\usetikzlibrary{shapes.geometric,calc,positioning,arrows.meta}
\pagestyle{empty}

% Angle convention used by all macros + .cells.tex fragments: 0° at 12
% o'clock, increasing clockwise. TikZ's native polar convention is
% counter-clockwise from 3 o'clock; the helper emits the *clockwise-from-12*
% angle directly, so a TikZ rotation of {-90} on each disc tikzpicture
% restores the visual convention. We do that at the wheel level, not here.

% \volvellecell{<bech32-char>}: presentation hook for one bottom-disc cell.
% The helper-emitted .cells.tex fragments wrap each char in this macro and
% provide the (angle:radius) placement themselves. Kept tiny; print-legible
% at typical scaling.
\newcommand{\volvellecell}[1]{\textsf{\fontsize{6}{7}\selectfont #1}}

% \volvellecellinner{<bech32-char>}: same as \volvellecell but ~60% of the
% regular font size. Used by the helper for row 0 only — the innermost ring
% has the smallest arc length per cell (radius shrinks linearly with row),
% so its 32 cells crowd together at full font size.
\newcommand{\volvellecellinner}[1]{\textsf{\fontsize{3.6}{4.2}\selectfont #1}}

% \volvelleendpoint{<angle-deg>}{<format-label>}: thick arc segment at the
% named angle on the top disc + small text label outside the disc. The arc
% spans one segment (11.25°) centered on the angle. (Name avoids the
% LaTeX-reserved `\end…` prefix.)
\newcommand{\volvelleendpoint}[2]{%
  \draw[line width=2pt] ([shift=(#1-5.625:3.05in)]0,0)
    arc[start angle=#1-5.625, end angle=#1+5.625, radius=3.05in];
  \node at (#1:3.30in) [font=\sffamily\fontsize{7}{8}\selectfont] {#2};%
}

% \registrationtarget: center cross-hair + 8 outer alignment ticks at 45°
% intervals on the bottom disc. Helps the user verify cut+rotation accuracy.
\newcommand{\registrationtarget}{%
  \draw[line width=0.4pt] (-0.15in,0) -- (0.15in,0);
  \draw[line width=0.4pt] (0,-0.15in) -- (0,0.15in);
  \draw[line width=0.3pt] (0,0) circle (0.05in);
  \foreach \a in {0,45,90,135,180,225,270,315} {%
    \draw[line width=0.5pt] (\a:3.65in) -- (\a:3.85in);%
  }%
}

% \selftestlegend{<format>}{<target-residue>}{<canonical-input>}: bottom-
% margin self-test block in monospace. The user rotates through the
% canonical input; the disc must land on the format's endpoint band.
\newcommand{\selftestlegend}[3]{%
  \begin{minipage}{7.0in}\ttfamily\fontsize{7}{9}\selectfont\raggedright%
  format:          #1\\
  target residue:  #2\\
  canonical input: #3\\
  Rotate through the canonical input. The disc must land on the\\
  format-specific endpoint band. If it doesn't, this wheel is misprinted.%
  \end{minipage}%
}
